% !TEX root = ../../main.tex
\setchapterimage{figure/chapter2封面图.jpg}
\setchapterintro{}
\chapter{线性方程组}

\section{线性方程组}
\begin{definition}\phantomsection\label{p1:concept:c2-l7}
	一个线性方程组，有以下的形式
	\begin{equation}\tag{1}
		\left\{
		\begin{aligned}
			a_{11}x_1+\cdots+a_{1n}x_n&=b_1,\\
			a_{21}x_1+\cdots+a_{2n}x_n&=b_2,\\
			&\vdots\\
			a_{m1}x_1+\cdots+a_{mn}x_n&=b_m.
		\end{aligned}
		\right.
	\end{equation}
	其中 $a_{ij},b_i\in\F$，$x_1,\cdots,x_n$ 称为未知量。若 $b_i=0\ (1\leq i\leq m)$，则称 $(1)$ 为齐次的。
	
	若 $c_1,\cdots,c_n\in\F$，使得
	\[
	\left\{
	\begin{aligned}
		a_{11}c_1+\cdots+a_{1n}c_n&=b_1,\\
		a_{21}c_1+\cdots+a_{2n}c_n&=b_2,\\
		&\vdots\\
		a_{m1}c_1+\cdots+a_{mn}c_n&=b_m,
	\end{aligned}
	\right.
	\]
	那么，称 $(x_1,\cdots,x_n)=(c_1,\cdots,c_n)\in\F^n$ 为 $(1)$ 的一个解。显然的，$(0,0,\cdots,0)\in\F^n$ 是齐次方程的一个解。
	
	对另一个方程组
	\begin{equation}\tag{2}
		\left\{
		\begin{aligned}
			a'_{11}x_1+\cdots+a'_{1n}x_n&=b'_1,\\
			&\vdots\\
			a'_{l1}x_1+\cdots+a'_{ln}x_n&=b'_l,
		\end{aligned}
		\right.
	\end{equation}
	若它与 $(1)$ 有相同的解，那么称 $(1)$ 与 $(2)$ 是等价的。
\end{definition}

\begin{definition}\phantomsection\label{p1:concept:c2-l47}
	对 $(1)$，初等行变换规定为以下三种：
	\begin{enumerate}
		\item 将某一行 $(i)$ 乘以 $c\in\F$ 且 $c\neq0$，作为新方程组的 $(i)$，记为 $c(i)$；
		\item 将某一行 $(i)$ 的 $c$ 倍加到另一行 $(j)$，作为新的一行，记为 $c(i)+(j)$；
		\item 将行 $(i)$ 与行 $(j)$ 互换，记为 $(i)\leftrightarrow(j)$。
	\end{enumerate}
	
\end{definition}
显然地，初等行变换前后的方程组等价。



\newpage
\begin{definition}\phantomsection\label{p1:concept:c2-l61}
	\[
	\left\{
	\begin{aligned}
		a_{11}x_1+\cdots+a_{1n}x_n&=b_1,\\
		&\vdots\\
		a_{m1}x_1+\cdots+a_{mn}x_n&=b_m
	\end{aligned}
	\right.
	\]
	称为阶梯形的，若
	\[
	a_{i1}=a_{i2}=\cdots=a_{i,k-1}=0\quad\land\quad a_{ik}\neq0
	\]
	时，有
	\[
	a_{i+1,1}=a_{i+1,2}=\cdots=a_{i+1,k}=0.
	\]
	对 $1\leq i\leq m-1$，也就是说，一个阶梯形方程组，有以下的形式
	\[
	\left\{
	\begin{array}{rcl}
		a_{11}x_1+\cdots+a_{1k_1}x_{k_1}+\cdots a_{1k_2}x_{k_2}+\cdots+a_{1n}x_n&=&b_1,\\
		a_{2k_1}x_{k_1}+\cdots+a_{2k_2}x_{k_2}+\cdots+a_{2n}x_n&=&b_2,\\
		a_{3k_2}x_{k_2}+\cdots+a_{3n}x_n&=&b_3,\\
		&\vdots&\\
		a_{rk_r}x_{k_r}+\cdots+a_{rn}x_n&=&b_r,\\
		0&=&b_{r+1},\\
		&\vdots&\\
		0&=&b_m.
	\end{array}
	\right.
	\]
	其中 $a_{11},a_{2k_1},a_{3k_2},\cdots,a_{rk_r}$ 称为主元。
	
	若 $b_{r+1},\cdots,b_m\neq0$，则该方程组无解。
	
	称一个线性方程组是严格阶梯形，若有以下形式
	\[
	\left\{
	\begin{array}{rcl}
		a_{11}x_1+a_{12}x_2+\cdots+a_{1r}x_r+\cdots+a_{1n}x_n&=&b_1,\\
		a_{22}x_2+\cdots+a_{2r}x_r+\cdots+a_{2n}x_n&=&b_2,\\
		&\vdots&\\
		a_{rr}x_r+\cdots+a_{rn}x_n&=&b_r,\\
		0&=&b_{r+1},\\
		0&=&0,\\
		&\vdots&\\
		0&=&0.
	\end{array}
	\right.
	\]
	
\end{definition}
对一个变换成阶梯形的方程组，我们可以解出相应的根。
\begin{enumerate}
	\item $r=n$，对最后一个方程，解出
	\[
	x_n=\frac{b_n}{a_{nn}}.
	\]
	$x_i$ 可以递归地由
	\[
	x_i=\frac{1}{a_{ii}}\left(b_i-a_{i,i+1}x_{i+1}-\cdots-a_{in}x_n\right),\quad i=n-1,n-2,\cdots,1.
	\]
	\item $r<n$，则
	\[
	x_i=\frac{1}{a_{ii}}\left(b_i-a_{i,i+1}x_{i+1}-\cdots-a_{in}x_n\right),\quad i=r,\cdots,1.
	\]
	\phantomsection\label{p1:concept:c2-l129}$x_{r+1},\cdots,x_n$ 称为自由变量。
\end{enumerate}
\begin{lemma}\label{任意线性方程组可以通过有限多步初等行变换变为一个如上的阶梯形方程组}
	任意线性方程组可以通过有限多步初等行变换变为一个如上的阶梯形方程组
	\[
	\left\{
	\begin{aligned}
		a_{11}x_1+a_{12}x_2+\cdots+a_{1n}x_n&=b_1,\\
		a_{21}x_1+a_{22}x_2+\cdots+a_{2n}x_n&=b_2,\\
		&\vdots\\
		a_{m1}x_1+a_{m2}x_2+\cdots+a_{mn}x_n&=b_m.
	\end{aligned}
	\right.
	\]
\end{lemma}
\begin{proof}
	对 $n$ 做归纳。
	
	$n=1$ 时，方程组为
	\[
	\left\{
	\begin{aligned}
		a_{11}x_1&=b_1,\\
		a_{21}x_1&=b_2,\\
		&\vdots\\
		a_{m1}x_1&=b_m.
	\end{aligned}
	\right.
	\]
	\begin{enumerate}
		\item 若 $a_{11}\neq0$，则
		\[
		\left\{
		\begin{aligned}
			x_1&=\frac{b_1}{a_{11}},\\
			0&=b_2-\frac{a_{21}}{a_{11}}b_1,\\
			&\vdots\\
			0&=b_m-\frac{a_{m1}}{a_{11}}b_1.
		\end{aligned}
		\right.
		\]
		\item 若 $a_{11}=0$ 且 $\exists i\ \text{使得}\ a_{i1}\neq0$，则先互换得
		\[
		\left\{
		\begin{aligned}
			a_{i1}x_1&=b_i,\\
			a_{11}x_1&=b_1,\\
			&\vdots\\
			a_{m1}x_1&=b_m,
		\end{aligned}
		\right.
		\]
		之后同 $1$。
		\item 若 $\forall i=1,\cdots,m$，$a_{i1}=0$，则平凡。
	\end{enumerate}
	
	对 $n$：
	\begin{enumerate}
		\item 若 $a_{11}\neq0$，用行变换使各方程 $x_1$ 的系数为 $0$。对变换后且除去第一行的方程组得到的 $n-1$ 阶方程组用归纳法，再与第一行拼接即可。
		\item 若 $a_{11}=0$，取 $a_{i1}\neq0$，互换即可。
	\end{enumerate}
\end{proof}

对一个 $n$ 元齐次线性方程组
\[
\left\{
\begin{aligned}
	a_{11}x_1+\cdots+a_{1n}x_n&=0,\\
	&\vdots\\
	a_{m1}x_1+\cdots+a_{mn}x_n&=0,
\end{aligned}
\right.
\]
若有非零解，就有无穷多解。

如果方程的个数小于未知数的个数，也即 $m<n$，那么它一定有非零解。

\begin{definition}\phantomsection\label{p1:concept:c2-l206}
	称主元都是 $1$，且主元上方的系数都是 $0$ 的阶梯形方程组为行简化阶梯方程组。
	
	显然地，任意线性方程组都可以经过有限多步初等行变换变为行简化阶梯方程组（RREF）。
\end{definition}

\begin{proposition}
	任意线性方程组等价于唯一的行简化阶梯形方程组。
\end{proposition}
这个证明是不平凡的\red{(待补充)}




\section{向量}

\begin{definition}\phantomsection\label{p1:concept:c2-l222}
	一列或一行由 $a_i\in\F$ 组成的 $n$ 元组称为向量，用小写的希腊字母表示，通常 $\alpha,\beta,\gamma,\cdots$ 表示列向量，$\alpha',\beta',\gamma',\cdots$ 表示行向量，零向量用 $0$ 来表示。
	
	所有 $n$ 维列向量构成的集合称为 $n$ 维列向量空间，记为 $\F^n$；所有 $m$ 维行向量构成的集合称为 $m$ 维行向量空间，记为 $\F^m$。
	
	向量的加法与数乘不言自明。
	
	
\end{definition}
对方程组
\[
\left\{
\begin{aligned}
	a_{11}x_1+\cdots+a_{1n}x_n&=b_1,\\
	&\vdots\\
	a_{m1}x_1+\cdots+a_{mn}x_n&=b_m,
\end{aligned}
\right.
\]
可以写成
\[
x_1\alpha_1+\cdots+x_n\alpha_n=\beta
\quad\text{或}\quad
\sum_{k=1}^n x_k\alpha_k=\beta,
\]
其中
\[
\alpha_1=
\begin{pmatrix}
	a_{11}\\
	\vdots\\
	a_{m1}
\end{pmatrix},\quad
\cdots,\quad
\alpha_n=
\begin{pmatrix}
	a_{1n}\\
	\vdots\\
	a_{mn}
\end{pmatrix},\quad
\beta=
\begin{pmatrix}
	b_1\\
	\vdots\\
	b_m
\end{pmatrix}.
\]
研究线性方程组的解，等价于研究向量方程的解。
\begin{definition}\phantomsection\label{p1:concept:c2-l270}
	对 $\alpha_1,\cdots,\alpha_n\in\F^m$，与 $\beta\in\F^m$，若存在 $k_1,\cdots,k_n\in\F$，使得
	\[
	k_1\alpha_1+\cdots+k_n\alpha_n=\beta,
	\]
	则称 $\beta$ 可由 $\alpha_1,\cdots,\alpha_n$ 线性表示。
	
	若存在不全为 $0$ 的 $k_1,\cdots,k_n\in\F$，使
	\[
	k_1\alpha_1+\cdots+k_n\alpha_n=0,
	\]
	则称 $\alpha_1,\cdots,\alpha_n$ 线性相关；反之，定义出线性无关。
\end{definition}

\begin{proposition}
	若 $\beta$ 可由 $\alpha_1,\cdots,\alpha_n$ 线性表示，则表示方式不唯一当且仅当有无穷多种表示方式。
\end{proposition}
\begin{proof}
	充分性显然。
	
	必要性：设
	\[
	\beta=k_1\alpha_1+\cdots+k_n\alpha_n
	=l_1\alpha_1+\cdots+l_n\alpha_n,
	\]
	且存在 $i$，使得 $k_i\neq l_i$，那么
	\[
	0=(k_1-l_1)\alpha_1+\cdots+(k_n-l_n)\alpha_n.
	\]
	对任意 $t\in\F$，有
	\[
	0=t(k_1-l_1)\alpha_1+\cdots+t(k_n-l_n)\alpha_n.
	\]
	于是
	\[
	\beta=\beta+0=
	\sum_{j=1}^n\bigl(k_j+t(k_j-l_j)\bigr)\alpha_j.
	\]
	故 $\beta$ 就有无穷多种表示方式。
\end{proof}

我们可以得出以下结论。对 $\F$ 上的线性方程组
\[
x_1\alpha_1+\cdots+x_n\alpha_n=\beta,
\]
有
\begin{enumerate}
	\item 若无解：$\beta$ 不能由 $\alpha_1,\cdots,\alpha_n$ 线性表示；
	\item 有唯一一解：$\beta$ 可由 $\alpha_1,\cdots,\alpha_n$ 线性表示，且表示方式唯一；
	\item 无穷多解：$\beta$ 可由 $\alpha_1,\cdots,\alpha_n$ 线性表示，且表示方式不唯一。
\end{enumerate}

实际上，$\beta$ 表示方式的唯一性，等价于 $0$ 表示方式的唯一性。

\begin{proposition}
	$\beta$ 可由 $\alpha_1,\cdots,\alpha_n\in\F^m$ 线性表示，则表示方式唯一，当且仅当 $0$ 由 $\alpha_1,\cdots,\alpha_n$ 线性表示方式唯一。
\end{proposition}
\begin{proof}
	必要性：若不然，存在 $k_1,\cdots,k_n\in\F$，且 $k_1,\cdots,k_n$ 不全为 $0$，使得
	\[
	k_1\alpha_1+\cdots+k_n\alpha_n=0.
	\]
	对 $\beta=\sum l_i\alpha_i$，有
	\[
	\beta=\beta+0=\sum (l_i+k_i)\alpha_i,
	\]
	矛盾。
	
	充分性：若不然，设
	\[
	\beta=\sum_{i=1}^n k_i\alpha_i=
	\sum_{i=1}^n l_i\alpha_i.
	\]
	那么
	\[
	0=\sum_{i=1}^n(k_i-l_i)\alpha_i.
	\]
	因 $0$ 的表示方式唯一，$k_i-l_i=0$，矛盾。
\end{proof}

放在线性方程组上看，我们有
\begin{corollary}
	$x_1\alpha_1+\cdots+x_n\alpha_n=\beta$ 有解，则有唯一一解，当且仅当
	$x_1\alpha_1+\cdots+x_n\alpha_n=0$ 仅有平凡解。
\end{corollary}

用线性相关（无关）的语言叙述如下：
\[
\begin{aligned}
	x_1\alpha_1+\cdots+x_n\alpha_n=\beta\text{ 有解，则有唯一一解}
	&\Longleftrightarrow \alpha_1,\cdots,\alpha_n\text{ 线性无关},\\
	x_1\alpha_1+\cdots+x_n\alpha_n=\beta\text{ 有无穷多解}
	&\Longleftrightarrow \alpha_1,\cdots,\alpha_n\text{ 线性相关}.
\end{aligned}
\]

考虑一个向量组与它的子组，我们有以下的命题。

\begin{proposition}
	对向量组 $(I):\alpha_1,\cdots,\alpha_n$，有
	\begin{enumerate}
		\item 若有一个子组 $\alpha_{i_1},\cdots,\alpha_{i_r}$ 线性相关，则 $(I)$ 相关；
		\item 若任意子组 $\alpha_{i_1},\cdots,\alpha_{i_r}$ 无关，则 $(I)$ 无关。
	\end{enumerate}
\end{proposition}

\begin{proposition}\label{2.2.7}
	$\alpha_1,\cdots,\alpha_n\in\F^m$，$n>m$，则 $\alpha_1,\cdots,\alpha_n$ 线性相关。
\end{proposition}
\begin{proof}
	利用未知数个数大于方程个数时齐次方程必有非零解这一点，显而易见。
\end{proof}

\begin{definition}\phantomsection\label{p1:concept:c2-l383}
	设 $\alpha_1,\cdots,\alpha_n\in\F^m$，其中
	\[
	\alpha_i=
	\begin{pmatrix}
		a_{1i}\\
		\vdots\\
		a_{mi}
	\end{pmatrix}.
	\]
	取 $j_1<\cdots<j_k$，对 $b_{j_1 i},\cdots,b_{j_k i}\in\F$，令
	\[
	\widetilde{\alpha}_i=
	\begin{pmatrix}
		a_{1i}\\
		\vdots\\
		b_{j_1 i}\\
		\vdots\\
		b_{j_k i}\\
		\vdots\\
		a_{mi}
	\end{pmatrix}\in\F^{m+k}.
	\]
	称 $\widetilde{\alpha}_1,\widetilde{\alpha}_2,\cdots,\widetilde{\alpha}_n$ 是 $\alpha_1,\alpha_2,\cdots,\alpha_n$ 的伸长组。
\end{definition}

\begin{proposition}
	若 $\alpha_1,\cdots,\alpha_n$ 无关，则任意伸长组 $\widetilde{\alpha}_1,\cdots,\widetilde{\alpha}_n$ 无关。
\end{proposition}
\begin{proof}
	考虑
	\begin{align}
		x_1\alpha_1+\cdots+x_n\alpha_n&=0,\tag{1}\\
		x_1\widetilde{\alpha}_1+\cdots+x_n\widetilde{\alpha}_n&=0.\tag{2}
	\end{align}
	$(2)$ 的解集 $\subset (1)$ 的解集，而 $(1)$ 的解集 $=\{0\}$，且 $(2)$ 齐次，故 $(2)$ 有且仅有零解。
\end{proof}

\section{向量组之间的关系}

\begin{lemma}
	$\alpha_1,\cdots,\alpha_n$ 线性相关的充要条件是至少有一个 $\alpha_i$ 可以由其余的向量线性表示。
\end{lemma}
\begin{proof}
	必要性：由 $\alpha_1,\cdots,\alpha_n$ 线性相关，存在不全为 $0$ 的
	$k_1,\cdots,k_n\in\F$，使得
	\[
	\sum_{j=1}^n k_j\alpha_j=0.
	\]
	由于 $k_1,\cdots,k_n$ 不全为 $0$，故存在 $i$，使得 $k_i\neq0$。于是
	\[
	\alpha_i=\sum_{j\neq i}\left(-\frac{k_j}{k_i}\right)\alpha_j.
	\]

	充分性：显然。
\end{proof}

\begin{lemma}
	若 $\alpha_1,\cdots,\alpha_n$ 线性无关，且
	$\alpha_1,\cdots,\alpha_n,\beta$ 线性相关，则 $\beta$ 可由
	$\alpha_1,\cdots,\alpha_n$ 线性表示。
\end{lemma}
\begin{proof}
	存在不全为 $0$ 的 $k_1,\cdots,k_n,k\in\F$，使得
	\[
	\sum_{i=1}^n k_i\alpha_i+k\beta=0.
	\]
	若 $k=0$，则有 $\sum k_i\alpha_i=0$，而 $\alpha_1,\cdots,\alpha_n$ 无关，矛盾。因此 $k\neq0$，故
	\[
	\beta=-\sum_{i=1}^n\frac{k_i}{k}\alpha_i.
	\]
\end{proof}

\begin{theorem}(基本引理)
	对向量组 $(I):\alpha_1,\cdots,\alpha_r$ 与
	$(II):\beta_1,\cdots,\beta_s$，若
	\begin{enumerate}
		\item 每个 $\alpha_i$ 都可以由 $(II)$ 线性表示；
		\item $r>s$，
	\end{enumerate}
	则 $(I)$ 线性相关。
\end{theorem}
这是命题 \ref{2.2.7} 的推广，其逆否命题也会经常用到：
	\[
	\left.
	\begin{array}{l}
		\text{每个 }\alpha_i\text{ 都可以由 }(II)\text{ 线性表示},\\
		(I)\text{ 线性无关}
	\end{array}
	\right\}
	\Longrightarrow r\leq s.
	\]
\begin{proof}
	设
	\[
	\sum_{i=1}^r x_i\alpha_i=0,\qquad x_i\in\F.
	\]
	由条件 $(1)$，对每个 $i$，存在 $k_{ij}\in\F$，使得
	\[
	\alpha_i=\sum_{j=1}^s k_{ij}\beta_j.
	\]
	于是
	\begin{align*}
		\sum_{i=1}^r x_i\sum_{j=1}^s k_{ij}\beta_j&=0,\\
		\sum_{j=1}^s\left(\sum_{i=1}^r k_{ij}x_i\right)\beta_j&=0.
	\end{align*}
	$(I)$ 线性相关的充分条件是上面系数
	$\sum_{i=1}^r k_{ij}x_i$ 全为 $0$。$k_{ij}$ 已知，因此只需找到不全为 $0$ 的 $x_i$，使得
	\[
	\sum_{i=1}^r k_{ij}x_i=0,\qquad 1\leq j\leq s,
	\]
	即齐次线性方程组
	\[
	\left\{
	\begin{aligned}
		k_{11}x_1+\cdots+k_{r1}x_r&=0,\\
		k_{12}x_1+\cdots+k_{r2}x_r&=0,\\
		&\vdots\\
		k_{1s}x_1+\cdots+k_{rs}x_r&=0.
	\end{aligned}
	\right.
	\]
	由于 $s<r$，由\ref{2.2.7}，该方程组有非平凡解 $x_1^0,\cdots,x_r^0$，也即
	\[
	\begin{pmatrix}
		x_1^0\\
		\vdots\\
		x_r^0
	\end{pmatrix}
	\neq0,
	\qquad
	\sum_{i=1}^r k_{ij}x_i^0=0,\quad 1\leq j\leq s.
	\]
	于是
	\[
	\sum_{i=1}^r x_i^0\alpha_i
	=\sum_{j=1}^s0\beta_j=0.
	\]
	故 $\alpha_1,\cdots,\alpha_r$ 线性相关。
\end{proof}

\begin{corollary}
	对上述 $(I)$、$(II)$，若 $(I)$、$(II)$ 均无关，且 $(I)$ 可以表示任意 $(II)$ 中的 $\beta_i$，$(II)$ 可以表示任意 $(I)$ 中的 $\alpha_j$，则 $r=s$。
\end{corollary}

\begin{definition}\phantomsection\label{p1:concept:c2-l528}
	$(I):\alpha_1,\cdots,\alpha_r$ 和 $(II):\beta_1,\cdots,\beta_s$，若 $(I)$ 中的每个向量都可以由 $(II)$ 线性表示，则称 $(I)$ 可由 $(II)$ 表出。显然，向量组的表出关系具有传递性。

	若 $(I)$ 与 $(II)$ 可以互相线性表示，则称它们是等价的，记为 $(I)\cong(II)$。显然，这是等价关系。

	显然，若 $(I)$ 是 $(II)$ 的子组，则 $(I)$ 可由 $(II)$ 线性表示。

	我们研究的
	\[
	x_1\alpha_1+\cdots+x_n\alpha_n=\beta
	\]
	有解的充要条件是
	\[
	\alpha_1,\cdots,\alpha_n
	\quad\text{与}\quad
	\alpha_1,\cdots,\alpha_n,\beta
	\]
	等价。
\end{definition}

\begin{definition}\phantomsection\label{p1:concept:c2-l548}
	对一个向量组 $(I):\alpha_1,\cdots,\alpha_n$，若其子组
	$(I'):\alpha_{i_1},\cdots,\alpha_{i_r}$ 满足
	\begin{enumerate}
		\item $\alpha_{i_1},\cdots,\alpha_{i_r}$ 线性无关；
		\item 任意 $\alpha_j$ 都可以由 $\alpha_{i_1},\cdots,\alpha_{i_r}$ 线性表示，
	\end{enumerate}
	则称 $(I')$ 为 $(I)$ 的极大无关子组。

	显然，不全为 $0$ 的向量所构成的向量组一定含有极大无关子组，但是不一定是唯一的。由基本引理，任意向量组的所有极大无关子组所包含的向量个数相同。

	称 $\alpha_1,\cdots,\alpha_n$ 的极大无关子组所含的向量个数为
	$\alpha_1,\cdots,\alpha_n$ 的秩，记为
	\[
	\operatorname{rank}(\alpha_1,\cdots,\alpha_n)
	\quad\text{或}\quad
	r(\alpha_1,\cdots,\alpha_n).
	\]
\end{definition}

\begin{proposition}
	一个向量组的秩具有以下性质：
	\begin{enumerate}
		\item 若 $\alpha_1,\cdots,\alpha_n\in\F^m$，则
		\[
		0\leq r(\alpha_1,\cdots,\alpha_n)\leq\min\{m,n\},
		\]
		且
		\[
		\alpha_1,\cdots,\alpha_n\text{ 线性相关}
		\Longleftrightarrow
		r(\alpha_1,\cdots,\alpha_n)<n.
		\]
		\item 若 $r(\alpha_1,\cdots,\alpha_n)=s$，那么组中的 $s+1$ 个向量必然线性相关。
		\item 若 $(I)$ 可由 $(II)$ 线性表出，则
		\[
		r(I)\leq r(II).
		\]
		\item 若 $(I)\cong(II)$，则
		\[
		r(I)=r(II).
		\]
	\end{enumerate}
\end{proposition}

\begin{proposition}
	若 $r(I)=r(II)$，且 $(I)$ 可以由 $(II)$ 表出，则 $(I)\cong(II)$。
\end{proposition}
\begin{proof}
	取 $(I)$ 的极大无关子组
	\[
	(I'):\alpha_{i_1},\cdots,\alpha_{i_r},
	\]
	取 $(II)$ 的极大无关子组
	\[
	(II'):\beta_{j_1},\cdots,\beta_{j_s}.
	\]
	由 $r(I)=r(II)$ 有 $r=s$。

	由于 $(I)\cong(I')$ 且 $(II)\cong(II')$，又 $(I)$ 可以由 $(II)$ 表出，故 $(I')$ 可以由 $(II')$ 表出。下证 $(II')$ 可以由 $(I')$ 表出。

	考虑
	\[
	(III):\alpha_{i_1},\cdots,\alpha_{i_r},\beta_{j_1},\cdots,\beta_{j_r}.
	\]
	那么 $(III)$ 的极大无关组是 $(II')$，所以 $r(III)=r$。断言 $(I')$ 也是 $(III)$ 的一个极大无关子组。

	若 $\beta_{j_k}$ 不能由 $(I')$ 表出，则
	\[
	\beta_{j_k},\alpha_{i_1},\cdots,\alpha_{i_r}
	\]
	线性无关，与 $r(III)=r$ 矛盾。这样 $(I')$ 也是 $(III)$ 的一个极大无关子组，得证。
\end{proof}

使用秩的语言，我们不难对
\[
x_1\alpha_1+\cdots+x_n\alpha_n=\beta
\]
解的情况进行叙述：
\begin{enumerate}
	\item 方程组有解，当且仅当
	\[
	r(\alpha_1,\cdots,\alpha_n)
	=
	r(\alpha_1,\cdots,\alpha_n,\beta).
	\]
	\item 有唯一一解，当且仅当
	\[
	r(\alpha_1,\cdots,\alpha_n)
	=
	r(\alpha_1,\cdots,\alpha_n,\beta)
	=n.
	\]
	\item 有无穷多解，当且仅当
	\[
	r(\alpha_1,\cdots,\alpha_n)
	=
	r(\alpha_1,\cdots,\alpha_n,\beta)
	<n.
	\]
\end{enumerate}



\section{矩阵与线性方程组}

\begin{definition}\phantomsection\label{p1:concept:c2-l654}
	称数表
	\[
	\begin{pmatrix}
		a_{11}&\cdots&a_{1n}\\
		\vdots&&\vdots\\
		a_{m1}&\cdots&a_{mn}
	\end{pmatrix}
	\]
	为 $\F$ 上的一个 $m\times n$（行数 $\times$ 列数）矩阵，$a_{ij}\in\F$ 称为它的 $(i,j)$--元。矩阵通常用大写字母 $A,B,\cdots$ 来表示，有时用 $A_{m\times n}$ 来标识其尺寸。
	
	规定
	\[
	A_{m\times n}=B_{m'\times n'}
	\]
	当且仅当 $m=m'$，$n=n'$，且
	\[
	\forall\,1\leq i\leq m,
	\quad 1\leq j\leq n,
	\quad A(i,j)=B(i,j).
	\]
	当 $m=n$ 时，我们称该矩阵为方阵。
	
	在 $(i,j)$--元的位置留空，则默认其为 $0$；写 $*$，则默认此处未知。
	用 $M_{m\times n}(\F)$ 或 $\F^{m\times n}$ 表示 $\F$ 上所有 $m\times n$ 型矩阵所构成的集合。特别地，当 $n=1$ 时，$M_{m\times 1}(\F)$ 等价于 $\F^m$，即 $m$ 维列向量空间。
	
	$A_{m\times n}$ 的第 $i$ 个列向量记为
	\[
	\alpha^i=
	\begin{pmatrix}
		A(1,i)\\
		\vdots\\
		A(m,i)
	\end{pmatrix}\in\F^m.
	\]
	第 $j$ 个行向量记为
	\[
	\beta_j=\bigl(A(j,1)\quad \cdots\quad A(j,n)\bigr)\in\F^n.
	\]
	那么 $A$ 可以表示为列向量形式或行向量形式
	\[
	A_{m\times n}=(\alpha^1\ \cdots\ \alpha^n)
	=
	\begin{pmatrix}
		\beta_1\\
		\vdots\\
		\beta_m
	\end{pmatrix}.
	\]
\end{definition}

\begin{definition}\phantomsection\label{p1:concept:c2-l705}
	对任意线性方程组
	\[
	\left\{
	\begin{aligned}
		a_{11}x_1+\cdots+a_{1n}x_n&=b_1,\\
		&\vdots\\
		a_{m1}x_1+\cdots+a_{mn}x_n&=b_m,
	\end{aligned}
	\right.
	\]
	定义其系数矩阵为 $A=(a_{ij})\in M_{m\times n}(\F)$，即
	\[
	A=
	\begin{pmatrix}
		a_{11}&a_{12}&\cdots&a_{1n}\\
		a_{21}&a_{22}&\cdots&a_{2n}\\
		\vdots&\vdots&\ddots&\vdots\\
		a_{m1}&a_{m2}&\cdots&a_{mn}
	\end{pmatrix}.
	\]
	定义其增广矩阵为
	\[
	A_{\mathrm{aug}}=
	\begin{pmatrix}
		a_{11}&a_{12}&\cdots&a_{1n}&b_1\\
		a_{21}&a_{22}&\cdots&a_{2n}&b_2\\
		\vdots&\vdots&\ddots&\vdots&\vdots\\
		a_{m1}&a_{m2}&\cdots&a_{mn}&b_m
	\end{pmatrix}
	\in M_{m\times(n+1)}(\F).
	\]
\end{definition}

\begin{definition}\phantomsection\label{p1:concept:c2-l739}
	矩阵 $A$ 的行（列）向量组的秩称为 $A$ 的行（列）秩。显然地，
	\[
	\forall A\in M_{m\times n}(\F),\qquad
	0\leq A\text{ 的行（列）秩}\leq \min\{m,n\}.
	\]
\end{definition}

\begin{definition}\phantomsection\label{p1:concept:c2-l747}
	若矩阵 $J$ 满足：
	\begin{enumerate}
		\item $J$ 的零行（如果有的话）都位于 $J$ 的下方；
		\item $J$ 的每个非零行的第一个非零元的列指标严格递增，
	\end{enumerate}
	则称 $J$ 是一个阶梯形矩阵。
	
	阶梯形矩阵可以写成
	\[
	J=
	\begin{pmatrix}
		0&\cdots&C_{1j_1}&\cdots&C_{1j_2}&\cdots&C_{1j_r}&\cdots&C_{1n}\\
		0&\cdots&0&\cdots&C_{2j_2}&\cdots&C_{2j_r}&\cdots&C_{2n}\\
		\vdots&&\vdots&&\vdots&&\vdots&&\vdots\\
		0&\cdots&0&\cdots&0&\cdots&C_{rj_r}&\cdots&C_{rn}\\
		0&\cdots&0&\cdots&0&\cdots&0&\cdots&0\\
		\vdots&&\vdots&&\vdots&&\vdots&&\vdots\\
		0&\cdots&0&\cdots&0&\cdots&0&\cdots&0
	\end{pmatrix},
	\]
	其中
	\[
	1\leq j_1<\cdots<j_r\leq n.
	\]
	默认 $C_{1j_1},\cdots,C_{rj_r}$ 均不为 $0$，称
	$C_{1j_1},\cdots,C_{rj_r}$ 为主元。
\end{definition}

\begin{theorem}
	对上述的 $J$，有行秩等于列秩等于 $r$。
\end{theorem}
\begin{proof}
	取 $J$ 的列向量为 $\alpha_1,\cdots,\alpha_n$，行向量为
	$\beta_1,\cdots,\beta_r,0,\cdots,0$。
	
	先看行秩。行秩为 $r(\beta_1,\cdots,\beta_r)$。下证
	$\beta_1,\cdots,\beta_r$ 无关。设
	\[
	x_1\beta_1+\cdots+x_r\beta_r=0.
	\]
	则有
	\[
	\left\{
	\begin{aligned}
		x_1C_{1j_1}&=0,\\
		x_1C_{1j_2}+x_2C_{2j_2}&=0,\\
		&\vdots\\
		x_1C_{1j_r}+x_2C_{2j_r}+\cdots+x_rC_{rj_r}&=0.
	\end{aligned}
	\right.
	\]
	故 $x_1=\cdots=x_r=0$，所以 $\beta_1,\cdots,\beta_r$ 无关。
	
	再看列秩。下证 $\alpha_{j_1},\cdots,\alpha_{j_r}$ 是 $J$ 的列向量组的一个极大无关子组。
	
	首先，$\alpha_{j_1},\cdots,\alpha_{j_r}$ 无关。设
	\[
	x_1\alpha_{j_1}+\cdots+x_r\alpha_{j_r}=0.
	\]
	则有
	\[
	\left\{
	\begin{aligned}
		x_rC_{rj_r}&=0,\\
		x_{r-1}C_{r-1,j_{r-1}}+x_rC_{r-1,j_r}&=0,\\
		&\vdots\\
		x_1C_{1j_1}+\cdots+x_rC_{1j_r}&=0.
	\end{aligned}
	\right.
	\]
	故 $x_r=\cdots=x_1=0$。
	
	其次，对任意 $\alpha_k$，可由 $\alpha_{j_1},\cdots,\alpha_{j_r}$ 线性表示。设
	\[
	\alpha_k=
	\begin{pmatrix}
		C_{1k}\\
		C_{2k}\\
		\vdots\\
		C_{rk}\\
		0\\
		\vdots\\
		0
	\end{pmatrix},
	\qquad
	\alpha_k=z_1\alpha_{j_1}+\cdots+z_r\alpha_{j_r}.
	\]
	等价于
	\[
	\left\{
	\begin{aligned}
		C_{1k}&=z_1C_{1j_1}+\cdots+z_rC_{1j_r},\\
		&\vdots\\
		C_{r-1,k}&=z_{r-1}C_{r-1,j_{r-1}}+z_rC_{r-1,j_r},\\
		C_{rk}&=z_rC_{rj_r}.
	\end{aligned}
	\right.
	\]
	回代可知 $z_1,\cdots,z_r$ 可解，故得证。
\end{proof}

显然地，我们有以下结论：若 $J$ 是 $\F$ 上任意阶梯形矩阵，则
\begin{enumerate}
	\item $J$ 的主元所在列构成了 $J$ 的列向量组的一个极大无关子组；
	\item $J$ 的列秩等于 $J$ 的行秩，等于 $J$ 的主元个数。
\end{enumerate}

类似于线性方程组上的变换，我们看矩阵上的初等变换。

\begin{definition}\phantomsection\label{p1:concept:c2-l857}
	对 $\F$ 上的矩阵，我们定义初等行（列）变换。称满足以下规则的变换为初等行（列）变换：
	\begin{enumerate}
		\item 交换 $i,j$ 行（列）的位置，记为 $(i,j)$；
		\item 把第 $i$ 行（列）的 $c$ 倍后加到第 $j$ 行（列），记为 $c(i)+(j)$；
		\item 用 $c\neq0$，$c\in\F$ 数乘第 $i$ 行（列），记为 $c(i)$。
	\end{enumerate}
	我们把所有的初等行变换与初等列变换统称为初等变换。
\end{definition}

类似于 \ref{任意线性方程组可以通过有限多步初等行变换变为一个如上的阶梯形方程组}，我们有以下结果。

\begin{lemma}
	$\forall A\in M_{m\times n}(\F)$，可以经过有限多步初等行变换，将之变为一个不唯一的阶梯形矩阵。
\end{lemma}

\phantomsection\label{p1:concept:c2-l873}对于主元为 $1$，主元上方元素均为 $0$ 的阶梯形矩阵，我们称之为行简化阶梯形矩阵。

$\forall A\in M_{m\times n}(\F)$，可以经过有限多步 $(i),(ii),(iii)$ 型初等行变换变为行简化阶梯形矩阵，并且唯一。

设
\[
E_r=
\begin{pmatrix}
	1&0&\cdots&0\\
	0&1&\cdots&0\\
	\vdots&\vdots&\ddots&\vdots\\
	0&0&\cdots&1
\end{pmatrix}
\]
为 $r$ 阶单位矩阵，形如
\[
\begin{pmatrix}
	E_r&0\\
	0&0
\end{pmatrix}
\]
\phantomsection\label{p1:concept:c2-l894}的矩阵，称为标准形矩阵。

$\forall A\in M_{m\times n}(\F)$，可以经过有限多步初等变换变为标准形矩阵。

对上面的变换 $(i),(ii),(iii)$，分别考虑变换前后的极大无关组，不难得出初等行（列）变换不改变行（列）秩。因此，我们得到以下结论。

\begin{theorem}
	$\forall A\in M_{m\times n}(\F)$，行秩等于列秩。
\end{theorem}

\begin{definition}\phantomsection\label{p1:concept:c2-l904}
	$A\in M_{m\times n}(\F)$，$A$ 的秩为 $A$ 的行秩或列秩，记为 $r(A)$ 或 $\operatorname{rank}(A)$，有
	\[
	r(A)\leq\min\{m,n\}.
	\]
	因域的扩张不改变线性相关或无关的关系，故矩阵的秩不依赖于域的扩张。
	
	
\end{definition}
对 $A,B\in M_{m\times n}(\F)$，若 $A$ 能够经过有限步初等变换为 $B$，就称 $A,B$ 相抵。

使用 $r(A)$ 与相抵的概念，我们对之前所述内容进行重述。

\begin{lemma}
	$A\in M_{m\times n}(\F)$，则 $r(A)=r$ 当且仅当 $A$ 相抵于
	\[
	\begin{pmatrix}
		E_r&0\\
		0&0
	\end{pmatrix},
	\]
	且与 $A$ 相抵的标准形矩阵唯一。
\end{lemma}
\begin{proof}
	显然。
\end{proof}

\begin{proposition}
	$A,B\in M_{m\times n}(\F)$，$A$ 与 $B$ 相抵当且仅当
	\[
	r(A)=r(B).
	\]
\end{proposition}
\begin{proof}
	必要性：由上一引理显然。
	
	充分性：由 $r(A)=r(B)=r$，则
	\[
	A\text{ 相抵于 }
	\begin{pmatrix}
		E_r&0\\
		0&0
	\end{pmatrix},
	\qquad
	B\text{ 相抵于 }
	\begin{pmatrix}
		E_r&0\\
		0&0
	\end{pmatrix},
	\]
	故 $A,B$ 相抵。
\end{proof}

利用矩阵的秩的语言，我们给出线性方程组解存在性定理的矩阵版本。

\begin{proposition}\label{2.4.12}
	对
	\[
	Ax=\beta,
	\qquad
	A_{\mathrm{aug}}=[A\ \beta],
	\]
	有
	\begin{enumerate}
		\item 有解，当且仅当
		\[
		r(A)=r(A_{\mathrm{aug}});
		\]
		\item 有唯一一解，当且仅当
		\[
		r(A)=r(A_{\mathrm{aug}})=n;
		\]
		\item 有无穷多解，当且仅当
		\[
		r(A)=r(A_{\mathrm{aug}})<n.
		\]
	\end{enumerate}
\end{proposition}

\begin{proposition}
	设 $\eta_1,\cdots,\eta_m$ 线性无关，$A=(a_{ij})_{m\times n}$，令
	\[
	\alpha_i=\sum_{j=1}^m a_{ji}\eta_j,
	\qquad 1\leq i\leq n.
	\]
	则
	\[
	r(\alpha_1,\cdots,\alpha_n)=r(A),
	\]
	且若 $\beta_{j_1},\cdots,\beta_{j_r}$ 为 $A$ 的列向量组的一个极大无关组，则
	$\alpha_{j_1},\cdots,\alpha_{j_r}$ 是 $\alpha_1,\cdots,\alpha_n$ 的一个极大无关组。
\end{proposition}
\begin{proof}
	不妨一般性，设 $A$ 的列向量组的一个极大无关组为
	$\beta_{j_1},\cdots,\beta_{j_r}$。下证
	$\alpha_{j_1},\cdots,\alpha_{j_r}$ 是 $\alpha_1,\cdots,\alpha_n$ 的极大无关组。
	
	$1.$ 先证 $\alpha_{j_1},\cdots,\alpha_{j_r}$ 无关。设
	\[
	x_1\alpha_{j_1}+\cdots+x_r\alpha_{j_r}=0.
	\]
	则
	\begin{align*}
	0
	&=\sum_{i=1}^r x_i\alpha_{j_i}\\
	&=\sum_{k=1}^r x_k\left(\sum_{t=1}^m a_{tj_k}\eta_t\right)\\
	&=\sum_{t=1}^m\left(\sum_{k=1}^r x_k a_{tj_k}\right)\eta_t.
	\end{align*}
	由 $\eta_1,\cdots,\eta_m$ 无关，得
	\[
	\sum_{k=1}^r x_k a_{tj_k}=0,
	\qquad 1\leq t\leq m.
	\]
	对
	\[
	\beta_{j_k}=
	\begin{pmatrix}
		a_{1j_k}\\
		\vdots\\
		a_{mj_k}
	\end{pmatrix},
	\]
	上式等价于
	\[
	\sum_{k=1}^r x_k\beta_{j_k}=0.
	\]
	而 $\beta_{j_1},\cdots,\beta_{j_r}$ 无关，故 $x_k=0$，$1\leq k\leq r$，从而
	$\alpha_{j_1},\cdots,\alpha_{j_r}$ 无关。
	
	$2.$ 再证任意 $\alpha_p$ 均由 $\alpha_{j_1},\cdots,\alpha_{j_r}$ 表出。由于
	$\beta_{j_1},\cdots,\beta_{j_r}$ 是 $A$ 的列向量组的一个极大无关组，对 $\beta_p$，存在 $l_1,\cdots,l_r\in\F$，使得
	\[
	\beta_p=\sum_{i=1}^r l_i\beta_{j_i}.
	\]
	从而 $\beta_p$ 的第 $j$ 个分量为
	\[
	a_{jp}=\sum_{i=1}^r l_i a_{j j_i},
	\qquad 1\leq j\leq m.
	\]
	于是
	\begin{align*}
	\alpha_p
	&=\sum_{j=1}^m a_{jp}\eta_j\\
	&=\sum_{j=1}^m\left(\sum_{i=1}^r l_i a_{j j_i}\right)\eta_j\\
	&=\sum_{i=1}^r l_i\left(\sum_{j=1}^m a_{j j_i}\eta_j\right)\\
	&=\sum_{i=1}^r l_i\alpha_{j_i}.
	\end{align*}
	得证。
\end{proof}


\section{线性方程组解的结构}

\begin{definition}\phantomsection\label{p1:concept:c2-l1057}
	$Ax=\beta$ 的导出组为 $Ax=0$。
	
	
\end{definition}
在 $Ax=\beta$ 有解的前提下，
	\[
	Ax=\beta \text{ 有唯一一解}
	\quad\Longleftrightarrow\quad
	Ax=0 \text{ 仅有平凡解}.
	\]
	反过来讲，
	\[
	Ax=\beta \text{ 有无穷多解}
	\quad\Longleftrightarrow\quad
	Ax=0 \text{ 有无穷多解}.
	\]
\begin{lemma}
	设 $(\mathrm{I})$：$Ax=\beta$ 是非齐次线性方程组，$(\mathrm{II})$ 是 $(\mathrm{I})$ 的导出组 $Ax=0$，则
	\begin{enumerate}
		\item $(\mathrm{I})$ 中的任意两解之差是 $(\mathrm{II})$ 的解；
		\item $(\mathrm{I})$ 的任意解与 $(\mathrm{II})$ 的任意解之和仍是 $(\mathrm{I})$ 的解。
	\end{enumerate}
\end{lemma}
\begin{proof}
	将 $(\mathrm{I})$、$(\mathrm{II})$ 分别写成
	\[
	(\mathrm{I})\quad \sum_{i=1}^n x_i\alpha_i=\beta,
	\qquad
	(\mathrm{II})\quad \sum_{i=1}^n x_i\alpha_i=0.
	\]
	设
	\[
	\begin{pmatrix}
		c_1\\
		\vdots\\
		c_n
	\end{pmatrix},
	\qquad
	\begin{pmatrix}
		d_1\\
		\vdots\\
		d_n
	\end{pmatrix}
	\]
	是 $(\mathrm{I})$ 的解，即
	\[
	\sum_{i=1}^n c_i\alpha_i=\beta,
	\qquad
	\sum_{i=1}^n d_i\alpha_i=\beta.
	\]
	那么
	\[
	\sum_{i=1}^n(c_i-d_i)\alpha_i=\beta-\beta=0,
	\]
	也即
	\[
	\begin{pmatrix}
		c_1-d_1\\
		\vdots\\
		c_n-d_n
	\end{pmatrix}
	\]
	是 $(\mathrm{II})$ 的解。
	
	第 $2$ 条同理可得。
\end{proof}

\begin{corollary}
	设 $Ax=\beta$ 的解集为 $S_1$，$Ax=0$ 的解集为 $S_2$。任取 $Ax=\beta$ 的一个解 $\delta$，则
	\[
	S_1=\{\delta+r\mid r\in S_2\}.
	\]
\end{corollary}
\begin{proof}
	$1.$ 先证
	\[
	S_1\supset \{\delta+r\mid r\in S_2\}.
	\]
	由上一引理第 $2$ 条可得。
	
	$2.$ 再证
	\[
	S_1\subset \{\delta+r\mid r\in S_2\}.
	\]
	任取 $\alpha\in S_1$，则
	\[
	\alpha=\delta+(\alpha-\delta).
	\]
	其中 $\delta\in S_1$，由上一引理第 $1$ 条可得 $\alpha-\delta\in S_2$，故得证。
\end{proof}

显然地，$Ax=0$ 的解集对解向量的线性组合封闭；对 $Ax=\beta$ 则不然。

\begin{definition}\phantomsection\label{p1:concept:c2-l1151}
	对 $Ax=0$，若有一个解向量组 $\xi_1,\cdots,\xi_s$ 满足：
	\begin{enumerate}
		\item $\xi_1,\cdots,\xi_s$ 线性无关；
		\item $Ax=0$ 的任意解可以由 $\xi_1,\cdots,\xi_s$ 表出，
	\end{enumerate}
	则称 $\xi_1,\cdots,\xi_s$ 为 $Ax=0$ 的基础解系。
\end{definition}

下面，我们证明基础解系的存在性。

\begin{theorem}
	对一个 $n$ 元齐次线性方程组 $Ax=0$，存在由 $n-r(A)$ 个解向量构成的基础解系。
\end{theorem}
\begin{proof}
	对
	\[
	\left\{
	\begin{aligned}
		a_{11}x_1+\cdots+a_{1n}x_n&=0,\\
		a_{21}x_1+\cdots+a_{2n}x_n&=0,\\
		&\vdots\\
		a_{m1}x_1+\cdots+a_{mn}x_n&=0,
	\end{aligned}
	\right.
	\]
	下面构造一个基础解系。
	
	对
	\[
	A=
	\begin{pmatrix}
		a_{11}&\cdots&a_{1n}\\
		\vdots&&\vdots\\
		a_{m1}&\cdots&a_{mn}
	\end{pmatrix}
	\]
	进行初等行变换，化为行简化阶梯形矩阵。并且，不失一般性地，不妨设主元列恰是前 $r=r(A)$ 列。那么可化为
	\[
	\begin{pmatrix}
		1&0&\cdots&0&C_{1,r+1}&\cdots&C_{1n}\\
		0&1&\cdots&0&C_{2,r+1}&\cdots&C_{2n}\\
		\vdots&\vdots&\ddots&\vdots&\vdots&&\vdots\\
		0&0&\cdots&1&C_{r,r+1}&\cdots&C_{rn}\\
		0&0&\cdots&0&0&\cdots&0\\
		\vdots&\vdots&&\vdots&\vdots&&\vdots\\
		0&0&\cdots&0&0&\cdots&0
	\end{pmatrix}.
	\]
	于是原方程组等价于
	\[
	\left\{
	\begin{aligned}
		x_1+C_{1,r+1}x_{r+1}+\cdots+C_{1n}x_n&=0,\\
		x_2+C_{2,r+1}x_{r+1}+\cdots+C_{2n}x_n&=0,\\
		&\vdots\\
		x_r+C_{r,r+1}x_{r+1}+\cdots+C_{rn}x_n&=0.
	\end{aligned}
	\right.
	\]
	我们有两类变量：
	\begin{enumerate}
		\item $x_1,\cdots,x_r$：主元变量；
		\item $x_{r+1},\cdots,x_n$：自由变量（共 $n-r$ 个）。
	\end{enumerate}
	那么一旦自由变量的值确定下来，主元变量的值亦被确定：
	\[
	x_i=-C_{i,r+1}x_{r+1}-\cdots-C_{in}x_n,
	\qquad 1\leq i\leq r.
	\]
	也即
	\begin{equation}\tag{1}
	\begin{pmatrix}
		x_1\\
		x_2\\
		\vdots\\
		x_r\\
		x_{r+1}\\
		\vdots\\
		x_n
	\end{pmatrix}
	=
	x_{r+1}
	\begin{pmatrix}
		-C_{1,r+1}\\
		-C_{2,r+1}\\
		\vdots\\
		-C_{r,r+1}\\
		1\\
		\vdots\\
		0
	\end{pmatrix}
	+\cdots+
	x_n
	\begin{pmatrix}
		-C_{1n}\\
		-C_{2n}\\
		\vdots\\
		-C_{rn}\\
		0\\
		\vdots\\
		1
	\end{pmatrix}.
	\end{equation}
	
	令 $t=n-r$，取 $t$ 个向量 $\eta_1,\cdots,\eta_t$ 如下：
	\[
	\eta_k=
	\begin{pmatrix}
		-C_{1,r+k}\\
		-C_{2,r+k}\\
		\vdots\\
		-C_{r,r+k}\\
		0\\
		\vdots\\
		1\\
		\vdots\\
		0
	\end{pmatrix},
	\qquad
	\text{第 } r+k \text{ 个位置为 }1,
	\qquad 1\leq k\leq t.
	\]
	下证 $\eta_1,\cdots,\eta_t$ 是基础解系。
	
	$1.$ 无关性。设
	\[
	s_1\eta_1+\cdots+s_t\eta_t=0.
	\]
	考虑第 $r+k$ 个分量，有
	\[
	s_k=0.
	\]
	故 $\forall i$，$s_i=0$，因此 $\eta_1,\cdots,\eta_t$ 无关。
	
	$2.$ 表出性。任取解
	\[
	x=\begin{pmatrix}x_1\\ \vdots\\ x_n\end{pmatrix},
	\]
	由 $(1)$ 知
	\[
	x=x_{r+1}\eta_1+\cdots+x_n\eta_t.
	\]
	得证。
\end{proof}

下面，考虑基础解系的性质。

\begin{proposition}
	对 $Ax=0$，有
	\begin{enumerate}
		\item 不同的基础解系等价，从而包含的解向量数量相同；
		\item 与一个基础解系等价的任意线性无关解向量组都构成一个基础解系；
		\item 设 $\xi_1,\cdots,\xi_{n-r}$ 是一个基础解系，那么解集可以表示为
		\[
		S=\left\{x\mid x=k_1\xi_1+\cdots+k_{n-r}\xi_{n-r},\ k_i\in\F\right\}.
		\]
	\end{enumerate}
\end{proposition}
\begin{proof}
	$1.$ 由基本引理得。
	
	$2.$ 由 $1$ 得。
	
	$3.$ 由上一存在性定理得。
\end{proof}

对 $Ax=0$，若 $r(A)<n$，则任意 $n-r(A)$ 个线性无关的解向量构成一个基础解系，且这不依赖于数域 $\F$ 的扩张。

综合我们以上讨论的内容，又能得出以下结论。

\begin{proposition}
	对 $Ax=\beta$，且
	\[
	r(A)=r(A_{\mathrm{aug}})<n.
	\]
	设 $\xi_1,\cdots,\xi_{n-r}$ 是 $Ax=0$ 的基础解系，任取 $Ax=\beta$ 的一个解 $\delta$。那么 $Ax=\beta$ 的任意解可以表示成
	\[
	x=\delta+k_1\xi_1+\cdots+k_{n-r}\xi_{n-r},
	\qquad k_i\in\F.
	\]
\end{proposition}

反过来想，对给定的一个向量集 $S$，是否一定有一个线性方程组，使 $S$ 是该线性方程组的解空间的基。于是有下面的命题。

\begin{proposition}
	令 $S$ 是一个由 $\F^n$ 中线性无关向量构成的向量集，那么存在一个 $n$ 个变量的线性齐次方程组，使得 $S$ 为这个方程组解的基础解系。
\end{proposition}

为了证明这个命题，我们需要下面的引理。

\begin{lemma}
	设 $A\in M_{m\times n}(\F)$，且 $r(A)=m\leq n$。若 $A$ 的前 $m$ 列无关，则可以只用初等行变换把 $A$ 变为如下形式：
	\[
	\begin{pmatrix}
		1&0&\cdots&0&b_{1,m+1}&\cdots&b_{1n}\\
		0&1&\cdots&0&b_{2,m+1}&\cdots&b_{2n}\\
		\vdots&\vdots&\ddots&\vdots&\vdots&&\vdots\\
		0&0&\cdots&1&b_{m,m+1}&\cdots&b_{mn}
	\end{pmatrix}.
	\]
\end{lemma}
\begin{proof}
	由 $A$ 的前 $m$ 列无关，第一列的元素不能全为 $0$。所以，若 $A(1,1)=0$，可以通过交换行，使 $A(1,1)\neq0$。于是
	\[
	A\longrightarrow
	\begin{pmatrix}
		1&*&\cdots&*\\
		0&*&\cdots&*\\
		\vdots&\vdots&&\vdots\\
		0&*&\cdots&*
	\end{pmatrix}
	=:B.
	\]
	同理，$B$ 的第 $2$ 列中 $B(2,2),\cdots,B(m,2)$ 不能全为 $0$，故
	\[
	B\longrightarrow
	\begin{pmatrix}
		1&*&*&\cdots&*\\
		0&1&*&\cdots&*\\
		\vdots&\vdots&\vdots&&\vdots\\
		0&0&*&\cdots&*
	\end{pmatrix}
	=:C.
	\]
	重复上述过程，可得
	\[
	C\longrightarrow
	\begin{pmatrix}
		1&*&\cdots&*&*&\cdots&*\\
		0&1&\cdots&*&*&\cdots&*\\
		\vdots&\vdots&\ddots&\vdots&\vdots&&\vdots\\
		0&0&\cdots&1&*&\cdots&*
	\end{pmatrix}
	\longrightarrow
	\begin{pmatrix}
		1&0&\cdots&0&*&\cdots&*\\
		0&1&\cdots&0&*&\cdots&*\\
		\vdots&\vdots&\ddots&\vdots&\vdots&&\vdots\\
		0&0&\cdots&1&*&\cdots&*
	\end{pmatrix}.
	\]
	得证。
\end{proof}

下面，给出上一命题的证明。

\begin{proof}[命题证明]
	设 $S=\{\alpha_1,\cdots,\alpha_s\}\subseteq\F^n$ 是无关的列向量集，用 $\alpha'_1,\cdots,\alpha'_s$ 表示相应的行向量。令
	\[
	A=
	\begin{pmatrix}
		\alpha'_1\\
		\alpha'_2\\
		\vdots\\
		\alpha'_s
	\end{pmatrix},
	\]
	则 $A$ 是 $s\times n$ 型矩阵，且 $r(A)=s$。
	
	不妨设 $A$ 的前 $s$ 列无关（否则，交换即可）。对 $A$ 只做初等行变换，可得
	\[
	\begin{pmatrix}
		1&0&\cdots&0&C_{1,s+1}&\cdots&C_{1n}\\
		0&1&\cdots&0&C_{2,s+1}&\cdots&C_{2n}\\
		\vdots&\vdots&\ddots&\vdots&\vdots&&\vdots\\
		0&0&\cdots&1&C_{s,s+1}&\cdots&C_{sn}
	\end{pmatrix}
	=
	\begin{pmatrix}
		\beta'_1\\
		\vdots\\
		\beta'_s
	\end{pmatrix}.
	\]
	$\alpha_1,\cdots,\alpha_s$ 与 $\beta_1,\cdots,\beta_s$ 等价，且 $\beta_1,\cdots,\beta_s$ 无关。构造一个齐次线性方程组，使 $\beta_1,\cdots,\beta_s$ 为其基础解系即可：
	\begin{equation}\tag{1}
	\left\{
	\begin{aligned}
		x_{s+1}&=C_{1,s+1}x_1+\cdots+C_{s,s+1}x_s,\\
		x_{s+2}&=C_{1,s+2}x_1+\cdots+C_{s,s+2}x_s,\\
		&\vdots\\
		x_n&=C_{1n}x_1+\cdots+C_{sn}x_s.
	\end{aligned}
	\right.
	\end{equation}
	从而 $\alpha_1,\cdots,\alpha_s$ 是 $(1)$ 的基础解系。
\end{proof}

\section{有价值的习题}

\begin{problem}
	令$S$ 是两个向量$\alpha$和$\beta$的集合，证明$S$是线性无关的当且仅当存在一个$k\in\F$使得
	\[
	\alpha=k\beta\qquad\text{或}\qquad\beta=k\alpha.
	\]
\end{problem}

\begin{proof}
	必要性 $(\Rightarrow)$：由 $\alpha,\beta$ 线性相关，存在不全为 $0$ 的 $k_1,k_2\in\F$，使得
	\[
	k_1\alpha+k_2\beta=0.
	\]
	若 $k_1\neq0$，则
	\[
	\alpha=-\frac{k_2}{k_1}\beta.
	\]
	取 $k=-\dfrac{k_2}{k_1}\in\F$，即有 $\alpha=k\beta$。若 $k_2\neq0$，则
	\[
	\beta=-\frac{k_1}{k_2}\alpha.
	\]
	取 $k=-\dfrac{k_1}{k_2}\in\F$，即有 $\beta=k\alpha$。
	
	充分性 $(\Leftarrow)$：若 $\alpha=k\beta$，则
	\[
	\alpha+(-k)\beta=0.
	\]
	这个线性组合的系数不全为 $0$，故 $\alpha,\beta$ 线性相关。$\beta=k\alpha$ 的情形同理。
\end{proof}
\begin{problem}
	令$r_1,r_2,r_3$ 分别是这三个向量组的秩
	\[
	\alpha_1,\alpha_2,\ldots,\alpha_s;
	\qquad
	\beta_1,\beta_2,\ldots,\beta_t;
	\qquad
	\alpha_1,\alpha_2,\ldots,\alpha_s,
	\beta_1,\beta_2,\ldots,\beta_t,
	\]
	那么
	\[
	\max(r_1,r_2)\leq r_3\leq r_1+r_2.
	\]
\end{problem}

\begin{proof}
	不妨设 $\alpha_i$ 不全为 $0$，$\beta_j$ 不全为 $0$。
	
	任取向量组
	\[
	(I):\alpha_1,\ldots,\alpha_s
	\]
	的一个极大无关子组
	\[
	(I'):\alpha_{i_1},\ldots,\alpha_{i_{r_1}};
	\]
	任取向量组
	\[
	(II):\beta_1,\ldots,\beta_t
	\]
	的一个极大无关子组
	\[
	(II'):\beta_{j_1},\ldots,\beta_{j_{r_2}}.
	\]
	于是，向量组
	\[
	(III):\alpha_1,\ldots,\alpha_s,\beta_1,\ldots,\beta_t
	\]
	等价于
	\[
	(III'):\alpha_{i_1},\ldots,\alpha_{i_{r_1}},
	\beta_{j_1},\ldots,\beta_{j_{r_2}}.
	\]
	由此即得
	\[
	r_3=\rank(III)=\rank(III')\leq r_1+r_2.
	\]
	又由于 $(I')$ 是 $(III')$ 的一个子组，所以
	\[
	r_1=\rank(I)=\rank(I')
	\leq\rank(III')=\rank(III)=r_3.
	\]
	同理可证 $r_2\leq r_3$，所以
	\[
	\max(r_1,r_2)\leq r_3.
	\]
\end{proof}



\begin{problem}
	令$A=(a_{ij})_{n\times n}$ 是一个 $n\times n$ 的实矩阵且
	\[
	a_{ii}>0,\qquad a_{ij}<0\quad(i\neq j),
	\]
	且使得
	\[
	\sum_{j=1}^n a_{ij}>0
	\qquad\text{对于所有}i.
	\]
	证明 $r(A)=n$.
\end{problem}

\begin{proof}
	假设 $\rank(A)<n$，则以 $A$ 为系数矩阵的齐次线性方程组有非零解，即存在不全为 $0$ 的 $k_1,\ldots,k_n\in\R$，使得
	\begin{equation}\tag{1}
		\sum_{j=1}^n a_{ij}k_j=0,
		\qquad 1\leq i\leq n.
	\end{equation}
	不妨设
	\[
	|k_1|=\max\{|k_1|,\ldots,|k_n|\}.
	\]
	显然 $|k_1|>0$。由 $a_{1j}<0$（$2\leq j\leq n$）和 $(1)$（取 $i=1$）得
	\[
	\begin{aligned}
		|a_{11}k_1|
		&=|-a_{12}k_2-\cdots-a_{1n}k_n|\\
		&\leq |a_{12}k_2|+\cdots+|a_{1n}k_n|\\
		&=-a_{12}|k_2|-\cdots-a_{1n}|k_n|\\
		&\leq(-a_{12}-\cdots-a_{1n})|k_1|.
	\end{aligned}
	\]
	由 $a_{11}>0$ 且 $|k_1|>0$，可得
	\[
	a_{11}\leq-a_{12}-\cdots-a_{1n},
	\]
	这与
	\[
	a_{11}+a_{12}+\cdots+a_{1n}>0
	\]
	矛盾。因此 $r(A)=n$。
\end{proof}

\begin{problem}
	令$A=(a_{ij})_{n\times n}$ 是一个 $n\times n$ 的实矩阵且
	\[
	a_{ii}>0,\qquad a_{ij}<0\quad(i\neq j),
	\]
	并且使得
	\[
	\sum_{j=1}^n a_{ij}=0
	\qquad\text{对于所有}i.
	\]
	证明 $r(A)=n-1$.
\end{problem}

\begin{proof}
	利用
	\[
	\sum_{j=1}^n a_{ij}=0,
	\qquad 1\leq i\leq n,
	\]
	即 $A$ 的每行元素之和为 $0$，把 $A$ 的前 $n-1$ 列加到最后一列。由
	\[
	A=
	\begin{pmatrix}
		a_{11}&a_{12}&\cdots&a_{1,n-1}&a_{1n}\\
		a_{21}&a_{22}&\cdots&a_{2,n-1}&a_{2n}\\
		\vdots&\vdots&\ddots&\vdots&\vdots\\
		a_{n-1,1}&a_{n-1,2}&\cdots&a_{n-1,n-1}&a_{n-1,n}\\
		a_{n1}&a_{n2}&\cdots&a_{n,n-1}&a_{nn}
	\end{pmatrix}
	\]
	得到
	\[
	B=
	\begin{pmatrix}
		a_{11}&a_{12}&\cdots&a_{1,n-1}&0\\
		a_{21}&a_{22}&\cdots&a_{2,n-1}&0\\
		\vdots&\vdots&\ddots&\vdots&\vdots\\
		a_{n-1,1}&a_{n-1,2}&\cdots&a_{n-1,n-1}&0\\
		a_{n1}&a_{n2}&\cdots&a_{n,n-1}&0
	\end{pmatrix}.
	\]
	因此
	\begin{equation}\tag{1}
		\rank(A)=\rank(B)\leq n-1,
	\end{equation}
	因为 $B$ 有零列。
	
	（利用矩阵元素的性质进行初等变换，在矩阵中得到更多的 $0$。）
	
	考虑 $B$ 的前 $n-1$ 行和前 $n-1$ 列组成的子矩阵
	\[
	C=
	\begin{pmatrix}
		a_{11}&a_{12}&\cdots&a_{1,n-1}\\
		a_{21}&a_{22}&\cdots&a_{2,n-1}\\
		\vdots&\vdots&\ddots&\vdots\\
		a_{n-1,1}&a_{n-1,2}&\cdots&a_{n-1,n-1}
	\end{pmatrix}.
	\]
	对于 $C$ 的每行，$1\leq i\leq n-1$，有
	\[
	\sum_{j=1}^{n-1}a_{ij}
	=\left(\sum_{j=1}^n a_{ij}\right)-a_{in}
	=0-a_{in}=-a_{in}>0.
	\]
	且 $C$ 的元素满足
	\[
	a_{ii}>0,
	\qquad
	a_{ij}<0\quad(1\leq i\neq j\leq n-1).
	\]
	所以 $C$ 满足上一题的条件，从而
	\[
	\rank(C)=n-1.
	\]
	但 $C$ 是 $A$ 的子矩阵，所以
	\begin{equation}\tag{2}
		\rank(A)\geq\rank(C)=n-1.
	\end{equation}
	由 $(1)$、$(2)$ 即得
	\[
	\rank(A)=n-1.
	\]
\end{proof}

\begin{proposition}
	若 $B$ 是 $A$ 的子矩阵，则
	\[
	\rank(B)\leq\rank(A).
	\]
\end{proposition}

\begin{proof}
	不妨设 $B\neq0$。设 $B$ 的行向量组的一个极大无关组为
	\[
	\alpha_{i_1},\ldots,\alpha_{i_r}.
	\]
	则 $A$ 的第 $i_1,\ldots,i_r$ 行分别是
	$\alpha_{i_1},\ldots,\alpha_{i_r}$ 的伸长向量，因而仍线性无关。所以 $A$ 的行秩不小于 $r$，即
	\[
	\rank(A)\geq r=\rank(B).
	\]
\end{proof}

\begin{problem}
	令$A$ 是以下方程组的系数矩阵
	\begin{equation}\tag{*}
		\left\{
		\begin{aligned}
			a_{11}x_1+a_{12}x_2+\cdots+a_{1n}x_n&=b_1,\\
			a_{21}x_1+a_{22}x_2+\cdots+a_{2n}x_n&=b_2,\\
			&\vdots\\
			a_{n1}x_1+a_{n2}x_2+\cdots+a_{nn}x_n&=b_n,
		\end{aligned}
		\right.
	\end{equation}
	且
	\[
	B=
	\begin{pmatrix}
		a_{11}&a_{12}&\cdots&a_{1n}&b_1\\
		a_{21}&a_{22}&\cdots&a_{2n}&b_2\\
		\vdots&\vdots&\ddots&\vdots&\vdots\\
		a_{n1}&a_{n2}&\cdots&a_{nn}&b_n\\
		b_1&b_2&\cdots&b_n&0
	\end{pmatrix}.
	\]
	若 $r(A)=r(B)$, 那么方程组 $(*)$ 有解。 反过来对吗？为什么？
\end{problem}

\begin{proof}
	设 $A_{\mathrm{aug}}$ 为 $(*)$ 的增广矩阵。由于 $A_{\mathrm{aug}}$ 是 $B$ 的子矩阵，所以
	\[
	r(A_{\mathrm{aug}})\leq r(B).
	\]
	
	由题设 $r(A)=r(B)$ 可知
	\[
	r(A_{\mathrm{aug}})\leq r(A).
	\]
	但 $A$ 是 $A_{\mathrm{aug}}$ 的子矩阵，故
	\[
	r(A_{\mathrm{aug}})\geq r(A).
	\]
	所以
	\[
	r(A_{\mathrm{aug}})=r(A),
	\]
	从而方程组 $(*)$ 有解。
	
	反之不对。方程组 $(*)$ 有解等价于
	\[
	r(A_{\mathrm{aug}})=r(A).
	\]
	虽然 $r(B)\geq r(A_{\mathrm{aug}})$，但是不能由此推出
	$r(A_{\mathrm{aug}})=r(B)$。换言之，$(*)$ 有解等价于列向量
	\[
	\beta=(b_1,b_2,\ldots,b_n)'
	\]
	可由 $A$ 的列向量组线性表出，但是行向量 $\beta'$ 未必能由 $A$ 的行向量组线性表出。
	
	例如，考虑线性方程组
	\[
	\left\{
	\begin{aligned}
		x_1&=1,\\
		x_3&=2,\\
		2x_3&=4.
	\end{aligned}
	\right.
	\]
	它显然有解，但是系数矩阵
	\[
	A=
	\begin{pmatrix}
		1&0&0\\
		0&0&1\\
		0&0&2
	\end{pmatrix}
	\]
	的秩为 $2$，而
	\[
	B=
	\begin{pmatrix}
		1&0&0&1\\
		0&0&1&2\\
		0&0&2&4\\
		1&2&4&0
	\end{pmatrix}
	\]
	的秩为 $3$。这里，$B$ 的最后一个行向量不能由 $B$ 的前三个行向量线性表出。
\end{proof}










































































